\documentclass[11pt, a4paper]{article}

\input{bfw-style.tex}

\title{\vspace*{-1.5cm}\Huge\bfseries\color{bfwPrimaryDark}Aufgabenheft Session~6\\[0.4em]
  \large\color{bfwInkSoft}\textnormal{Wirtschaftskreislauf \& Marktsituationen --- Übungen im IHK-Klausurformat}}
\author{\textsf{IT Lernfeld 1 --- Das Unternehmen und die eigene Rolle im Betrieb beschreiben}}
\date{Selbstlernmaterial}

\begin{document}
\maketitle
\thispagestyle{empty}

\begin{merksatz}[title={\bfseries So nutzen Sie dieses Aufgabenheft}]
Bearbeiten Sie alle Aufgaben zunächst \emph{ohne} in die Lösungen zu schauen. Schreiben
Sie Ihre Antworten in vollständigen Sätzen --- die IHK-Prüfung verlangt formulierte
Begründungen, nicht bloße Stichworte. Beachten Sie die \emph{Operatoren} (nennen,
erläutern, beschreiben, begründen, beurteilen) --- sie geben den geforderten
Antwort-Tiefegrad vor. Erst danach öffnen Sie den Lösungsteil ab Seite~\pageref{loesungen}
und vergleichen.
\end{merksatz}

\bigskip

\textbf{Kontext:} Sie sind Auszubildende bei der \textbf{JIKU IT-Solutions GmbH} in Hamburg.
Ihre Ausbilderin Frau Meyerhoff erklärt Ihnen, dass JIKU in verschiedenen Märkten aktiv ist
--- als Anbieter von IT-Dienstleistungen gegenüber Unternehmenskunden und als Nachfrager
auf dem Beschaffungsmarkt für Hardware und Software.

\bigskip

\noindent
\begin{tabular}{|l|l|l|l|l|l|l|l|}
\hline
\textbf{Aufgabe} & 1 & 2 & 3 & 4 & 5 & 6 & \textbf{Gesamt}\\
\hline
\textbf{Punkte} & /10 & /8 & /10 & /10 & /8 & /10 & /56\\
\hline
\end{tabular}

\tableofcontents
\vspace{1em}
\hrule

\section{Aufgaben}

\subsection*{Aufgabe~1 --- Einfacher Wirtschaftskreislauf (10 Punkte)}

\begin{enumerate}[label=(\alph*)]
  \item \textbf{Beschreiben} Sie die zwei Wirtschaftssubjekte im einfachen Wirtschaftskreislauf
  und erläutern Sie, welche Ströme zwischen ihnen verlaufen. \hfill (4~P.)

  \vspace{6cm}

  \item \textbf{Erklären} Sie, warum Geld- und Güterströme im Wirtschaftskreislauf stets
  entgegengesetzt verlaufen, und geben Sie je ein konkretes Beispiel für einen Geldstrom
  und einen Güterstrom zwischen Haushalten und Unternehmen. \hfill (3~P.)

  \vspace{5cm}

  \item \textbf{Ordnen} Sie die folgenden Vorgänge dem richtigen Strom im einfachen
  Wirtschaftskreislauf zu. \hfill (3~P.)

  \bigskip

  \noindent
  \begin{tabular}{|p{6cm}|p{7cm}|}
  \hline
  \textbf{Vorgang} & \textbf{Strom}\\
  \hline
  Frau Müller kauft bei JIKU einen Laptop für privat & \\[0.8cm]
  \hline
  Ein JIKU-Techniker arbeitet 40 Stunden pro Woche & \\[0.8cm]
  \hline
  JIKU überweist Frau Meyerhoff ihr Monatsgehalt & \\[0.8cm]
  \hline
  \end{tabular}
\end{enumerate}

\subsection*{Aufgabe~2 --- Erweiterter Wirtschaftskreislauf (8 Punkte)}

\begin{enumerate}[label=(\alph*)]
  \item \textbf{Nennen} Sie die fünf Wirtschaftssubjekte des erweiterten Wirtschaftskreislaufs
  und erläutern Sie kurz die spezifische Rolle des Staates. \hfill (4~P.)

  \vspace{6cm}

  \item \textbf{Ordnen} Sie die folgenden Geldströme dem richtigen Wirtschaftssubjekt als
  Quelle und Ziel zu. \hfill (4~P.)

  \bigskip

  \noindent
  \begin{tabular}{|p{5.5cm}|p{2.5cm}|p{2.5cm}|}
  \hline
  \textbf{Geldstrom} & \textbf{Von} & \textbf{An}\\
  \hline
  JIKU erhält Investitionszulage für neue Server & & \\[0.8cm]
  \hline
  Herr Mayer zahlt Kfz-Steuer & & \\[0.8cm]
  \hline
  JIKU verkauft IT-Systeme nach Frankreich & & \\[0.8cm]
  \hline
  Bank schreibt Herrn Mayer Zinsen gut & & \\[0.8cm]
  \hline
  \end{tabular}
\end{enumerate}

\subsection*{Aufgabe~3 --- Marktarten und Marktformen (10 Punkte)}

\begin{enumerate}[label=(\alph*)]
  \item \textbf{Unterscheiden} Sie freien Markt und geschlossenen Markt sowie Gütermarkt
  und Faktorenmarkt. Geben Sie je ein Beispiel aus der IT-Branche. \hfill (4~P.)

  \vspace{6cm}

  \item \textbf{Ordnen} Sie die folgenden IT-Marktsituationen der richtigen Marktform zu
  und \textbf{begründen} Sie kurz. \hfill (6~P.)

  \bigskip

  \noindent
  \begin{tabular}{|p{6.5cm}|p{2.5cm}|p{3.5cm}|}
  \hline
  \textbf{Situation} & \textbf{Marktform} & \textbf{Begründung}\\
  \hline
  Microsoft Windows hält mehr als 70\,\% Marktanteil bei Desktop-Betriebssystemen & & \\[1.2cm]
  \hline
  AWS, Azure und Google Cloud teilen den Cloud-Markt unter sich auf & & \\[1.2cm]
  \hline
  Hunderte WordPress-Plugin-Anbieter und Millionen Nutzer & & \\[1.2cm]
  \hline
  Spezialisierter KI-Chip: 1 Hersteller, 2 staatliche Abnehmer & & \\[1.2cm]
  \hline
  \end{tabular}
\end{enumerate}

\subsection*{Aufgabe~4 --- Preisbildung am Markt (10 Punkte)}

Auf dem IT-Komponentenmarkt für einen bestimmten Prozessortyp gelten folgende Daten:

\bigskip

\noindent
\begin{tabular}{|l|l|l|}
\hline
\textbf{Preis in €} & \textbf{Angebotsmenge (Stück)} & \textbf{Nachfragemenge (Stück)}\\
\hline
300 & 200 & 1\,400\\
\hline
350 & 330 & 1\,100\\
\hline
400 & 550 & 950\\
\hline
450 & 650 & 650\\
\hline
500 & 750 & 400\\
\hline
550 & 900 & 250\\
\hline
\end{tabular}

\begin{enumerate}[label=(\alph*)]
  \item \textbf{Bestimmen} Sie den Gleichgewichtspreis und die Gleichgewichtsmenge. \hfill (2~P.)

  \vspace{3cm}

  \item \textbf{Beurteilen} Sie, ob bei einem Preis von 300\,€ ein Angebots- oder
  Nachfrageüberhang vorliegt, und berechnen Sie dessen Größe. \hfill (3~P.)

  \vspace{4cm}

  \item \textbf{Erklären} Sie, ob bei einem Marktpreis von 300\,€ ein Käufer- oder
  Verkäufermarkt vorliegt, und begründen Sie, in welche Richtung sich der Preis bewegt. \hfill (3~P.)

  \vspace{4cm}

  \item \textbf{Erläutern} Sie, wie sich eine Steuerentlastung der JIKU-Kunden auf die
  Nachfragekurve und den Gleichgewichtspreis auswirkt. \hfill (2~P.)

  \vspace{3.5cm}
\end{enumerate}

\subsection*{Aufgabe~5 --- Vollkommener Markt (8 Punkte)}

\begin{enumerate}[label=(\alph*)]
  \item \textbf{Nennen} Sie die fünf Bedingungen des vollkommenen Marktes und erläutern
  Sie jede kurz. \hfill (5~P.)

  \vspace{7cm}

  \item \textbf{Beurteilen} Sie, ob der Markt für IT-Dienstleistungen in Hamburg als
  vollkommener oder unvollkommener Markt einzustufen ist. \textbf{Begründen} Sie anhand
  von mindestens zwei Kriterien. \hfill (3~P.)

  \vspace{5cm}
\end{enumerate}

\subsection*{Aufgabe~6 --- Fallaufgabe: JIKU im Marktumfeld (10 Punkte)}

JIKU IT-Solutions analysiert sein Marktumfeld, um die Vertriebsstrategie für das nächste
Geschäftsjahr zu planen. Das Unternehmen bedient hauptsächlich mittelständische
Unternehmenskunden (B2B) und hat festgestellt, dass die Nachfrage nach IT-Infrastruktur
in Krisenzeiten stark zurückgeht.

\begin{enumerate}[label=(\alph*)]
  \item \textbf{Benennen} Sie den Transaktionsbereich, in dem JIKU überwiegend tätig ist,
  und erläutern Sie den Unterschied zum B2C-Bereich. \hfill (2~P.)

  \vspace{4cm}

  \item \textbf{Beschreiben} Sie je einen konjunkturellen, saisonalen und strukturellen
  Faktor, der die Nachfrage nach JIKU-Leistungen beeinflusst. Geben Sie je ein konkretes
  Beispiel. \hfill (6~P.)

  \vspace{7cm}

  \item \textbf{Erklären} Sie, welche Maßnahme JIKU ergreifen könnte, um bei einem
  Verkäufermarkt (z.\,B.\ Chip-Engpässe) die gestiegenen Einkaufspreise weiterzugeben
  oder zu vermeiden. \hfill (2~P.)

  \vspace{4cm}
\end{enumerate}

% =====================================================================
\newpage
\section*{Lösungshinweise für Lehrende}
\label{loesungen}
\addcontentsline{toc}{section}{Lösungshinweise für Lehrende}
% =====================================================================

\begin{merksatz}[title={\bfseries Hinweis zur Nutzung}]
Dieser Abschnitt ist ausschließlich für Lehrende bestimmt. Die Lösungen geben
Musterlösungen und typische Fehlerquellen an. Teilnehmende erhalten diesen Teil
erst nach der eigenständigen Bearbeitung.
\end{merksatz}

\subsection*{Lösung Aufgabe~1}

\begin{enumerate}[label=(\alph*)]
  \item Wirtschaftssubjekte: \textbf{Haushalte} und \textbf{Unternehmen}.
  Haushalte stellen Faktorleistungen (Arbeit, Boden, Kapital) bereit; Unternehmen zahlen
  dafür Einkommen (Löhne, Gehälter, Zinsen, Miete). Unternehmen produzieren Konsumgüter;
  Haushalte zahlen dafür Konsumausgaben.

  \item Güter- und Geldströme verlaufen entgegengesetzt, weil jeder Austausch eine
  Leistung \emph{und} eine Gegenleistung umfasst. Beispiel Geldstrom: JIKU zahlt Gehalt
  an Frau Meyerhoff. Beispiel Güterstrom: Frau Meyerhoff stellt ihre Arbeitskraft bereit.

  \item
  \begin{itemize}
    \item Frau Müller kauft Laptop: \textbf{Konsumausgabe} (Geldstrom von Haushalt zu Unternehmen)
    \item JIKU-Techniker arbeitet: \textbf{Faktorleistung} (Arbeit von Haushalt zu Unternehmen)
    \item JIKU überweist Gehalt: \textbf{Einkommen} (Geldstrom von Unternehmen zu Haushalt)
  \end{itemize}
\end{enumerate}

\subsection*{Lösung Aufgabe~2}

\begin{enumerate}[label=(\alph*)]
  \item Fünf Wirtschaftssubjekte des erweiterten Wirtschaftskreislaufs:
  \begin{itemize}
    \item \textbf{Haushalte:} Stellen Faktorleistungen bereit und konsumieren Güter.
    \item \textbf{Unternehmen:} Produzieren Güter und beschäftigen Arbeitskräfte.
    \item \textbf{Staat:} Vereinnahmt direkte Steuern von Haushalten und indirekte Steuern über Unternehmen; gibt diese als Staatsausgaben, Transferleistungen und Subventionen aus.
    \item \textbf{Banken:} Vermitteln zwischen Sparanlagen von Haushalten und Kreditvergabe an Unternehmen.
    \item \textbf{Ausland:} Nimmt an Exporten und Importen teil; beeinflusst den Wirtschaftskreislauf durch Handelsströme.
  \end{itemize}

  \textit{Typischer Fehler: Direkte und indirekte Steuern verwechseln. Direkte Steuern = Haushalt
  → Staat. Indirekte Steuern = Haushalt zahlt an Unternehmen, Unternehmen an Staat.}

  \item
  \begin{itemize}
    \item Investitionszulage: Von \textbf{Staat} an \textbf{Unternehmen} (Subvention)
    \item Kfz-Steuer: Von \textbf{Haushalt} an \textbf{Staat} (direkte Steuer)
    \item JIKU verkauft nach Frankreich: Von \textbf{Ausland} an \textbf{Unternehmen} (Exporterlös)
    \item Zinsgutschrift: Von \textbf{Bank} an \textbf{Haushalt} (Zinszahlung auf Ersparnisse)
  \end{itemize}
\end{enumerate}

\subsection*{Lösung Aufgabe~3}

\begin{enumerate}[label=(\alph*)]
  \item \textbf{Freier Markt:} Jeder hat Zugang (z.\,B.\ offener IT-Komponentenmarkt).
  \textbf{Geschlossener Markt:} Zugang beschränkt (z.\,B.\ Großhandelsmarkt für IT-Distributoren).
  \textbf{Gütermarkt:} Handel mit Waren/Dienstleistungen (z.\,B.\ Softwaremarkt).
  \textbf{Faktorenmarkt:} Handel mit Produktionsfaktoren (z.\,B.\ IT-Fachkräftearbeitsmarkt,
  Serverpark-Mietmarkt).

  \item
  \begin{itemize}
    \item Microsoft Windows > 70\,\%: \textbf{Angebotsmonopol / Quasimonopol} ---
    ein Anbieter dominiert, viele Nachfrager; ab 40\,\% gilt Marktbeherrschung (§\,18 Abs.\,4 GWB).
    \item AWS, Azure, Google: \textbf{Angebotsoligopol} --- wenige Anbieter, viele Nachfrager.
    \item Hunderte Plugin-Anbieter, Millionen Nutzer: \textbf{Polypol} --- viele Anbieter
    und viele Nachfrager; Marktpreis bildet sich frei.
    \item 1 Hersteller, 2 staatliche Abnehmer: \textbf{Beschränktes Angebotsmonopol} ---
    ein Anbieter, wenige (hier: zwei) Nachfrager; klassische Nischenmarktsituation.
  \end{itemize}
\end{enumerate}

\subsection*{Lösung Aufgabe~4}

\begin{enumerate}[label=(\alph*)]
  \item \textbf{Gleichgewichtspreis:} 450\,€; \textbf{Gleichgewichtsmenge:} 650 Stück
  (Angebot = Nachfrage = 650 Stück).

  \item Bei 300\,€: Angebot = 200 Stück, Nachfrage = 1\,400 Stück.
  Nachfrage > Angebot → \textbf{Nachfrageüberhang} von 1\,400 $-$ 200 = \textbf{1.200 Stück}.

  \item Bei Nachfrageüberhang übersteigt die Nachfrage das Angebot → \textbf{Verkäufermarkt}:
  Die Verkäufer können den Preis erhöhen, bis der Gleichgewichtspreis von 450\,€ erreicht ist.

  \item Einkommenserhöhung → mehr Kaufbereitschaft → \textbf{Nachfragekurve verschiebt sich nach
  rechts} (von N0 nach N1) → neuer Gleichgewichtspreis liegt \emph{höher} als zuvor,
  da bei jedem Preis nun mehr nachgefragt wird.
\end{enumerate}

\subsection*{Lösung Aufgabe~5}

\begin{enumerate}[label=(\alph*)]
  \item Fünf Bedingungen des vollkommenen Marktes:
  \begin{itemize}
    \item \textbf{Vollkommene Konkurrenz:} Viele Anbieter und Nachfrager.
    \item \textbf{Güterhomogenität:} Alle Güter sind gleichartig.
    \item \textbf{Markttransparenz:} Alle Teilnehmer kennen alle Preise und Anbieter.
    \item \textbf{Präferenzlosigkeit:} Keine Markenpräferenz oder persönlichen Vorlieben.
    \item \textbf{Anpassungsfähigkeit:} Sofortige Reaktion auf Preisänderungen.
  \end{itemize}

  \item IT-Dienstleistungsmarkt Hamburg: \textbf{unvollkommener Markt}, weil mindestens
  zwei Bedingungen nicht erfüllt sind: (1)~\textbf{Keine Güterhomogenität} --- IT-Dienstleistungen
  unterscheiden sich stark in Qualität, Spezialisierung und Service. (2)~\textbf{Keine
  Markttransparenz} --- Kunden haben keinen vollständigen Überblick über alle Anbieter
  und Preise (Informationsasymmetrie).
\end{enumerate}

\subsection*{Lösung Aufgabe~6}

\begin{enumerate}[label=(\alph*)]
  \item JIKU ist überwiegend im \textbf{B2B-Bereich} tätig (Business-to-Business): Unternehmen
  handelt mit Unternehmen. Unterschied zu B2C (Business-to-Consumer): Im B2C-Bereich beliefert
  das Unternehmen Privatverbraucher (z.\,B.\ über einen Online-Shop); der B2B-Markt hat häufig
  höhere Auftragsvolumina, andere rechtliche Anforderungen und längere Vertragslaufzeiten.

  \item
  \begin{itemize}
    \item \textbf{Konjunkturell:} In einer Rezession kürzen Unternehmen IT-Investitionen;
    Nachfrage nach JIKU-Dienstleistungen sinkt (z.\,B.\ Stornierungen von Netzwerkprojekten).
    \item \textbf{Saisonal:} Gegen Jahresende erhöhen viele Unternehmen ihre IT-Ausgaben
    (Budgetreste), während im Sommer die Nachfrage schwächer ist.
    \item \textbf{Strukturell:} Die Umstellung auf Cloud-Computing verändert die Nachfrage
    nach klassischer On-Premise-Hardware zugunsten von Cloud-Dienstleistungen und -Beratung.
  \end{itemize}

  \item Bei Chip-Engpässen (Verkäufermarkt) kann JIKU (a) die Mehrkosten über \textbf{höhere
  Angebotspreise} an Kunden weitergeben, (b) durch \textbf{Rahmenverträge mit Lieferanten}
  Preissteigerungen abfedern oder (c) auf \textbf{Alternativkomponenten} ausweichen, um
  die Abhängigkeit von einem knappen Produkt zu reduzieren.
\end{enumerate}

\end{document}
